\section{Conclusion}\label{s:conclusion}

Understanding the modern state of language interoperability is a complex task. Everything from the lowest level of the programming stack to the highest level of abstractions plays a role in well languages play with oneanother. Despite this complexity there is a structure to it that can be understood, and furthermore utilizied to better understand how interoperability tools have evolved and likely will evolve. The goal of this paper was to build a framework for understanding the evolution of more traditionally based interoperability tools, and to build on this framework - and an existing tool - to explore more complex, nieche and possible future applications and asses their broader viability. The primary insight of this paper is that FFI based tools despite their long history still have considerable relevance and room to evolve further in the future. A secondary insight is that FFI and more unified execution models and their associated interoperability tools are not entirely mutually exclusive, which is to say that alot the more recent theoretical foundation of unified execution models can be applied to FFI based tools to improve their performance and usability. Ultimately its hard to gauge how things will evolve in the future, but its quite evident that language interoperability is a complex and multi-faceted problem deeply tied into the evolution of programming languages.   